\scanpage{64}

\unsure{kummérien|kumérien} \inserted{connexe} (ramifié) de \inserted{$S$, de} groupe de Galois $\mu_a$, le diagramme \deleted{étant} induit par le changement de base $\eta\longrightarrow S$ étant cartésien. Un tel diagramme est unique à isomorphisme près, \deleted{$X'$ est connexe,} et $f':X'\longrightarrow S'$ satisfait aux mêmes hypothèses \deleted{ici} que $f$ dans \qquad, mais l'entier $d(X'/S')$ (défini comme $d=d(X/S)$ dans (2)) est maintenant \deleted{une puissance} égal\deleted{\unsure{x}} à $p^r$.
Si $I'\subset I$ est le groupe d'inertie \deleted{\unsure{de}} relatif à $S'$, alors pour tout entier $n$ premier à $p$, \unsure{à}l'homomorphisme
\[
u'_n:H^1(X',\mu_n)\longrightarrow H^1(X_{\bar\eta},\mu_n)^{I'}
\]
est injectif, et son conoyau est isomorphe à $\deleted{\unsure{\Gamma}}{}_n\delta/{}_n\delta'$, où $\delta$ est le sous-groupe de torsion de $\operatorname{Coker}w$ ($w$ défini dans (5)) et $\delta'$ le sous-groupe de torsion de $\operatorname{NS}(X_0)$. On en conclut, pour tout nombre premier $\ell\neq p$, \deleted{une} que l'homomorphisme
\[
u'(\ell):H^1(X',\underline Z_\ell\deleted{(1)})\longrightarrow H^1(X'_{\bar\eta},\underline Z_\ell\deleted{(1)})^{I'}
\]
est un \uline{isomorphisme}.\deleted{\unsure{et la même conclusion}}

\margin{$\times\quad\times\quad\times$}
\margin{63}
